Implementation direction: Context-Aware Design Invariants \section{Deterministic Scalable Design System}

Beyond structural and reactive semantics, a coherent UI system
requires deterministic rules for color contrast and spatial layout.
We formalize a Scalable Design System composed of two orthogonal
subsystems:

\begin{itemize}
    \item A discrete contrast-preserving color index model.
    \item A line-height–quantized spatial metric.
\end{itemize}

Both subsystems are finite, index-based, and independent of
component hierarchy.

\subsection{Indexed Contrast Color Model}

Let $T = \{0,1,\dots,11\}$ be a totally ordered finite tone index set.
Each element $t \in T$ represents a perceptual position between
background and contrast extremes.

For a given semantic color $c$, define a tone mapping:

\[
\Phi_c : T \rightarrow \mathcal{C}
\]

where $\mathcal{C}$ is the color space (e.g., sRGB).
The mapping is constructed such that:

\begin{equation}
|i - j| \geq 6 \quad \Rightarrow \quad
\text{Contrast}(\Phi_c(i), \Phi_c(j)) \geq 4.5:1
\end{equation}

This establishes a discrete contrast guarantee:
any pair of tones separated by at least six indices
satisfies WCAG AA accessibility for normal text.

\paragraph{Semantic Assignment.}

UI roles are expressed as deterministic index assignments:

\begin{itemize}
    \item Background surfaces: $t \in \{0,1,2,3,4,5\}$
    \item Primary text: $t = 8$
    \item Emphasis text: $t = 10$
    \item Secondary text: $t = 6$
    \item Borders: $t \in \{2,4,6\}$
    \item Elevation shadows: $t \in \{3,5,7\}$
\end{itemize}

Because indices are discrete and ordered,
semantic roles reduce to integer assignments rather than
independent color definitions.

\subsection{State Transitions as Index Jumps}

Interactive states are modeled as bounded index transitions.
Let $t$ be the base tone of an element. A state transformation is:

\[
\Psi(t, \Delta) = \text{clamp}(t + \Delta, 0, 11)
\]

where $\Delta$ is a fixed integer offset.

Examples:

\begin{itemize}
    \item Hover background: $\Delta = +4$
    \item Selected background: $\Delta = +2$
    \item Disabled background: $\Delta = -3$
    \item Focus ring: $\Delta = \pm 3$
\end{itemize}

Since all transformations operate in index space,
state styling becomes deterministic and contrast-preserving.
No additional color definitions are required.

\paragraph{Invariant.}

Let $d = |t_1 - t_2|$ denote tone distance.
Under any bounded state transition $\Delta$ with
$|\Delta| \leq 4$, relative contrast ordering is preserved:

\[
d' = |(t_1 + \Delta_1) - (t_2 + \Delta_2)|
\]

remains bounded and predictable.
Thus accessibility properties are stable under interaction.

\subsection{Line-Height Quantized Spacing System}

Let $L$ denote the typographic line height of the root context.
Define the base spatial quantum:

\[
u = \frac{L}{6}
\]

All layout dimensions are expressed as integer multiples of $u$:

\[
\mathcal{D} = \{ n u \mid n \in \mathbb{N} \}
\]

This forms a one-dimensional metric lattice over the vertical axis.

\paragraph{Rationale.}

Empirically, common text sizes produce:

\[
L \approx 1.5 \cdot \text{fontSize}
\]

Thus:

\[
u \approx 0.25\,\text{em}
\]

which corresponds to approximately $3$–$4$ pixels
for standard body text sizes.

This quantum allows fine-grained control of:

\begin{itemize}
    \item Vertical padding
    \item Border radius
    \item Component gaps
    \item Elevation offsets
\end{itemize}

while maintaining consistent vertical rhythm.

\subsection{Geometric Constraints}

To preserve compositional coherence, we enforce:

\begin{align}
\text{padding}_x &\geq \text{padding}_y \\
\text{borderRadius} &= \text{padding}_y \\
\text{gap} &\in [\text{padding}_y, \text{padding}_x]
\end{align}

Since all values belong to $\mathcal{D}$,
layout relationships remain algebraically consistent.

\subsection{Deterministic Design Layer}

The indexed color model and the quantized spacing system
share three structural properties:

\begin{enumerate}
    \item Finite ordered index domains.
    \item Deterministic arithmetic transformations.
    \item Independence from component hierarchy.
\end{enumerate}

Therefore, visual semantics are reduced to integer operations
over bounded sets.

This establishes a deterministic design layer
that complements the CARM execution model:
while CARM guarantees constant-time reactive updates,
the scalable design system guarantees bounded,
predictable visual transformations.

Together, they form a structurally flat,
computationally efficient, and perceptually coherent
architecture for declarative user interfaces, where tone and size are inherited from parent context and remain mutually consistent.
